\chapter{Cas Pratique : Entreprise LiZbiZ}
\label{chap:lizbiz}

\begin{boxlizbiz}
\textbf{Mission :} Vous êtes recrutés en tant que consultants externes par la Direction des Systèmes d'Information (DSI) de l'entreprise \textbf{LiZbiZ}. Votre objectif sur la durée du cours : sécuriser l'entreprise, auditer les pratiques et gérer les crises inhérentes à son activité.
\end{boxlizbiz}

% ------------------------------------------------------------------------------
\section{Identité de l'Entreprise}
% ------------------------------------------------------------------------------

\subsection{Présentation Générale}
\textbf{LiZbiZ} est une entreprise de e-commerce spécialisée dans la vente de matériel informatique et high-tech. Elle se distingue par sa promesse de livraison en 24h grâce à une logistique automatisée ("Smart Logistics").

\begin{itemize}
    \item \textbf{Secteur :} E-commerce \& Logistique.
    \item \textbf{Siège Social :} Lyon, France.
    \item \textbf{Effectif :} 150 salariés (100 logisticiens, 30 administratifs/commerciaux, 20 techniciens/ingénieurs).
    \item \textbf{Chiffre d'affaires :} 45 Millions d'euros.
    \item \textbf{Statut :} Société Anonyme (SA), introduite en bourse il y a 6 mois.
\end{itemize}

\subsection{Organigramme Simplifié}
La gouvernance est structurée ainsi :
\begin{itemize}
    \item \textbf{PDG (Jean-Pierre L.) :} Très médiatique, attaché à l'image de marque.
    \item \textbf{DSI (Sophie M.) :} Votre commanditaire. Souhaite moderniser la sécurité.
    \item \textbf{RSSI (Poste Vacant / Interim) :} Le poste est vacant, d'où votre recrutement.
    \item \textbf{DAF (Directeur Financier) :} Très sensible aux coûts et aux risques financiers.
    \item \textbf{Responsable Logistique :} Gère l'entrepôt connecté (IoT).
\end{itemize}

% ------------------------------------------------------------------------------
\section{Le Système d'Information (SI)}
% ------------------------------------------------------------------------------

Le SI de LiZbiZ est le cœur de l'entreprise. Il gère à la fois la vente en ligne et l'automatisation de l'entrepôt.

\subsection{Architecture Matérielle et Logicielle}

Le tableau ci-dessous récapitule les actifs matériels et logiciels critiques.

\begin{table}[htbp]
\centering
\begin{tabular}{p{3cm}p{4cm}p{6cm}}
    \toprule
    \textbf{Catégorie} & \textbf{Actif} & \textbf{Description Technique} \\
    \midrule
    \textbf{Infrastructure} & Serveur Principal & Dell PowerEdge R740 (ESXi) \\
                            & Baie de Stockage & NAS Synology (Données \& Sauvegardes) \\
                            & Réseau & Switch Cisco Catalyst (Core), Switch Netgear (Accès) \\
    \midrule
    \textbf{Logiciels} & ERP (Progiciel) & Odoo v15 (Gestion stocks, ventes, comptabilité) \\
                       & E-commerce & Site Web sous Magento 2 \\
                       & Base de Données & PostgreSQL 13 (Données clients \& commandes) \\
                       & Annuaire & Windows Server 2019 (Active Directory) \\
    \midrule
    \textbf{Périphériques} & Terminaux & 40 Postes Windows 10, 10 Laptops (Télétravail) \\
                           & IoT Logistique & 50 Scanners RFID (Android), 20 Tablettes de stock \\
    \bottomrule
\end{tabular}
\caption{Inventaire des actifs principaux de LiZbiZ}
\end{table}

\subsection{Topologie Réseau et Segmentation}

L'architecture réseau est segmentée pour séparer les flux. C'est cette topologie que vous retrouverez dans le laboratoire GNS3.

\begin{figure}[htbp]
\centering
% Note: Insérer ici le schéma réel généré sous GNS3 ou un schéma TikZ
% \includegraphics[width=0.8\textwidth]{figures/schema_reseau_lizbiz.png}
\begin{tikzpicture}[
    node distance=1.5cm,
    server/.style={rectangle, draw=cyberblue, thick, fill=cyberblue!10, minimum width=2.5cm, minimum height=1cm, align=center},
    firewall/.style={rectangle, draw=cyberred, thick, fill=cyberred!10, minimum width=2cm, minimum height=1cm, align=center},
    cloud/.style={ellipse, draw=black!50, thick, fill=gray!10, minimum width=2cm, minimum height=1cm}
]
    % Nodes
    \node[cloud] (internet) {Internet};
    \node[firewall, below=of internet] (fw) {Firewall\\(pFsense)};
    
    \node[server, below left=1.5cm and 2cm of fw] (dmz) {DMZ\\(Web/DNS)};
    \node[server, below=1.5cm of fw] (app) {APP\\(ERP/DB)};
    \node[server, below right=1.5cm and 2cm of fw] (user) {Users\\(AD/Postes)};
    \node[server, below=0.5cm of app] (iot) {IoT\\(Logistique)};

    % Connections
    \draw[->, thick] (internet) -- (fw);
    \draw[->, thick] (fw) -- (dmz);
    \draw[->, thick] (fw) -- (app);
    \draw[->, thick] (fw) -- (user);
    \draw[->, thick] (app) -- (iot);

\end{tikzpicture}
\caption{Schéma logique simplifié du réseau LiZbiZ}
\end{figure}

\textbf{Détail des VLANs :}
\begin{itemize}
    \item \textbf{VLAN 10 (DMZ) :} Accessible depuis Internet. Héberge le site marchand. Haut risque.
    \item \textbf{VLAN 20 (Applicatif) :} Cœur du système. Contient l'ERP et la Base de Données. Aucun accès direct externe autorisé.
    \item \textbf{VLAN 30 (Bureau/Utilisateurs) :} Postes de travail du siège. Accès Internet et Applicatif.
    \item \textbf{VLAN 40 (IoT/Logistique) :} Réseau isolé pour les scanners et capteurs. Faible sécurité par conception.
    \item \textbf{VLAN 99 (Administration) :} Accès de gestion aux équipements. Très sensible.
\end{itemize}

% ------------------------------------------------------------------------------
\section{Environnement de Laboratoire (Immersion)}
% ------------------------------------------------------------------------------

Pour simuler cette entreprise, nous utiliserons une maquette virtuelle.

\subsection{Configuration du Lab (GNS3 / EVE-NG)}
Les étudiants disposeront d'une image pré-configurée contenant :
\begin{itemize}
    \item \textbf{1 Routeur/Firewall pFsense :} Filtre les flux entre VLANs.
    \item \textbf{2 Machines Virtuelles (Linux Alpine/Debian) :} Simulant le serveur Web et l'ERP.
    \item \textbf{1 Serveur Windows Server :} Simulant l'Active Directory.
    \item \textbf{1 Attaquant (Kali Linux) :} Simulant les menaces externes (accès internet).
\end{itemize}

\subsection{Données Sensibles Simulées}
Pour réalisme, des jeux de données fictifs sont injectés :
\begin{itemize}
    \item \textbf{Clients :} 10 000 entrées (Nom, Prénom, Email, Adresse, Hash de mot de passe).
    \item \textbf{Transactions :} Historique de commandes et moyens de paiement (fictifs).
\end{itemize}

% ------------------------------------------------------------------------------
\section{Contexte Réglementaire et Contractuel}
% ------------------------------------------------------------------------------

LiZbiZ est soumise à des contraintes légales strictes, point de départ de l'analyse de risque.

\begin{enumerate}
    \item \textbf{RGPD :} Gestion de données personnelles clients et employés.
    \item \textbf{PCI-DSS :} Acceptation des paiements par carte bancaire (niveau 1).
    \item \textbf{Loi de Programmation Militaire (LPM) :} Opérateur d'importance vitale (OIV) potentiel car lié à la logistique critique (à débattre).
\end{enumerate}

% ------------------------------------------------------------------------------
\section{Problématiques Actuelles}
% ------------------------------------------------------------------------------

\begin{boxalert}
\textbf{Situation initiale :} Suite à une croissance rapide, la sécurité a été négligée. Le PDG s'inquiète d'une récente vague d'attaques par Ransomware chez des concurrents. La DSI manque de visibilité sur son niveau de risque réel.
\end{boxalert}

\textbf{Points de douleur identifiés :}
\begin{itemize}
    \item Mises à jour de sécurité irrégulières sur le serveur Web.
    \item Absence de plan de continuité d'activité (PCA) formalisé.
    \item Droits d'accès trop larges sur l'ERP (anciens collaborateurs encore actifs).
\end{itemize}

C'est sur cette base que commence votre mission.


% ==============================================================================
% Fichier : 02_module_risque.tex (Extrait Chapitre 1)
% Description : Introduction à la Gouvernance de la Sécurité
% ==============================================================================